\question[3]
\begin{parts}
\part Schreibe eine rekursive Funktion \texttt{quersumme(n)}, die die Quersumme einer natürlichen Zahl
zurückgibt. Tipp: \texttt{n \% 10} ist die letzte Ziffer, \texttt{n // 10} die Zahl ohne die letzte Ziffer.
\part Gib die Abfolge der rekursiven Aufrufe bei \texttt{quersumme(4711)} an.
\end{parts}
\begin{solutionbox}{5cm}
\begin{lstlisting}
# a)
def quersumme(n):
    if n == 0: return 0
    return n % 10 + quersumme(n // 10)
\end{lstlisting}
b) \texttt{quersumme(4711)}, \texttt{quersumme(471)}, \texttt{quersumme(47)}, \texttt{quersumme(4)},
\texttt{quersumme(0)}, Ergebnis $1 + 1 + 7 + 4 + 0 = 13$
\end{solutionbox}
